\chapter{Preparing the Pi}
\label{ch:pisetup}

\chapterbrief
  {Install Java and Maven, enable SPI, join the hardware groups, add swap.}
  {Why a 512\,MB machine needs care.}
  {A Pi that can build and run this project.}

\section{The fast way}

The starter pack ships a script that does everything in this chapter. Read the
chapter anyway --- when something fails later you want to know what was done.

\begin{shell}
# on your laptop, from the template folder
scp scripts/pi-bootstrap.sh brewer:~
ssh brewer 'bash pi-bootstrap.sh'
\end{shell}

Then reboot and skip to the checkpoint:

\begin{shell}
ssh brewer 'sudo reboot'
\end{shell}

The rest of this chapter is the same work, done by hand.

\section{Update first}

\begin{shell}
# on the Pi
sudo apt update
sudo apt full-upgrade -y
\end{shell}

On a Zero 2 W this takes a while on a fresh image. Let it finish.

\begin{concept}{\texttt{update} versus \texttt{upgrade}}
\code{apt update} refreshes the \emph{catalogue} --- what versions exist. It
installs nothing.

\code{apt upgrade} installs newer versions of what you already have.
\code{full-upgrade} additionally allows removing packages when a dependency
requires it, which is what you want on a fresh system.

Running \code{upgrade} without \code{update} first is the most common apt
mistake: it dutifully upgrades to the newest versions it knows about, which may
be months old.
\end{concept}

\section{Java and Maven}

\begin{shell}
sudo apt install -y openjdk-21-jdk maven git rsync
java -version
mvn -v
\end{shell}

\begin{trace}
openjdk version "21.0.11" 2026-04-21
OpenJDK Runtime Environment (build 21.0.11+10-1-deb13u2-Debian)
\end{trace}

Maven on the Pi is used rarely --- Chapter~\ref{ch:deploy} builds on your laptop
--- but it is needed once, to download the dependency jars.

\section{Enable SPI}

SPI is the bus the NFC reader speaks. It is off by default.

\begin{shell}
sudo raspi-config
\end{shell}

Navigate: \textbf{3 Interface Options} $\rightarrow$ \textbf{I4 SPI}
$\rightarrow$ \textbf{Yes} $\rightarrow$ \textbf{Ok} $\rightarrow$
\textbf{Finish}.

Or without the menu:

\begin{shell}
sudo raspi-config nonint do_spi 0      # 0 means ENABLE, confusingly
\end{shell}

\section{Join the hardware groups}

Linux does not let arbitrary programs poke hardware. Access is granted by group
membership.

\begin{shell}
sudo usermod -aG spi,gpio $USER
\end{shell}

\begin{trap}{Group changes need a fresh login}
Your current session still has your \emph{old} groups. Adding yourself to a group
does nothing until you log out and back in --- and over SSH that means
disconnecting, not just opening a new shell.

This produces one of the most confusing errors in the whole setup: you ran the
command, you can see the group in \code{/etc/group}, and the program still says
\code{Permission denied}.

\begin{shell}
exit
ssh brewer
groups          # must now list: spi gpio
\end{shell}
\end{trap}

\section{Swap, on a 512\,MB machine}

A Zero 2 W has 512\,MB of RAM. A JVM plus Maven will exceed that, and when Linux
runs out of memory the kernel starts killing processes --- frequently
\code{sshd}, which locks you out until you power-cycle it.

\begin{shell}
free -h                    # check what you have

sudo dphys-swapfile swapoff
sudo nano /etc/dphys-swapfile
\end{shell}

Change \code{CONF\_SWAPSIZE=100} to \code{CONF\_SWAPSIZE=1024}, then:

\begin{shell}
sudo dphys-swapfile setup
sudo dphys-swapfile swapon
free -h                    # Swap: should now show ~1.0Gi
\end{shell}

\begin{concept}{Swap is a safety net, not speed}
Swap is disk space used as overflow memory. Disk is thousands of times slower
than RAM, so a machine actively swapping is a machine crawling.

That is fine --- the goal is not performance, it is \emph{not dying}. Without
swap, running out of memory kills your SSH session. With swap, it merely gets
slow, and you can still type.

On an SD card, heavy swapping does wear the card. Which is one more argument for
building on your laptop.
\end{concept}

\section{Reboot and verify}

\begin{shell}
sudo reboot
\end{shell}

Wait thirty seconds, reconnect, and check everything:

\begin{shell}
ssh brewer

ls -l /dev/spidev0.0        # the SPI device must exist
groups                      # must include spi and gpio
java -version               # 21.x
free -h                     # swap around 1 GB
uname -m                    # aarch64
\end{shell}

\begin{trace}
crw-rw---- 1 root spi 153, 0 Jul 26 14:02 /dev/spidev0.0
martin adm dialout cdrom sudo audio video plugdev games users input
       render netdev spi gpio lpadmin
\end{trace}

Read that first line: the device is owned by group \code{spi}, and the group has
read/write (\code{rw}). That is exactly why you needed the group membership.

\section{One optional nicety}

Give yourself a clock and a tidy history:

\begin{shell}
# on the Pi
cat >> ~/.bashrc <<'EOF'

# timestamps in history -- surprisingly useful when debugging
export HISTTIMEFORMAT="%F %T  "
export HISTSIZE=5000
export HISTCONTROL=ignoredups

# shortcuts
alias ll='ls -alh'
alias ..='cd ..'
EOF

source ~/.bashrc
\end{shell}

\begin{checkpoint}{The Pi is ready}
\begin{shell}
ssh brewer 'ls -l /dev/spidev0.0 && groups | tr " " "\n" | grep -E "^(spi|gpio)$" && java -version'
\end{shell}

You should see the SPI device, both groups, and Java 21.

If \code{/dev/spidev0.0} does not exist, SPI is not enabled --- rerun
\code{raspi-config} and \emph{reboot}. If the groups are missing, log out and
back in.
\end{checkpoint}
